% Provenance (never printed): the markets-with-algebra strand of ../economics_series.md (grade 9: the percentage
% change in quantity for a 1 % change in price; elastic and inelastic; revenue when the price rises; necessities and
% luxuries; short run and long run). Brief: sources/book1/BRIEFS.md, g9-02. Point and arc elasticities are grade 10's.
% Chart data: figdata/grade-9/02-elasticity-first-look (published elasticities, each a ledger row; BLS gasoline average
% prices; an invented demand). Ledger: g9:andreyeva-2010-* (16), g9:fuel-elasticity-short, g9:fuel-elasticity-long,
% g8:cigarette-tax-elasticity, g9:eia-gasoline-2022, g9:south-africa-cigarettes.
\chapter{Elasticity: A First Look}\label{ch:g9:elasticity-first-look}

\begin{openingchart}{How much buyers cut their purchases when the price rises by $1\,\%$: published
estimates, in percent. Foods: averages of United States studies, 1938--2007. Cigarettes: rich
countries. Gasoline: averages of studies from many countries, in the first year (short run) and after
several years (long run).}{Andreyeva, Long and Brownell (2010), Table 1; National Cancer Institute and
World Health Organization (2016); Havranek, Irsova and Janda (2012)}
\begin{tikzpicture}
\begin{axis}[omrecord, xbar, bar width=6pt, y=0.36cm, xmin=0, xmax=1,
  ytick=data, yticklabels from table={figdata/grade-9/02-elasticity-first-look/out/elasticities.csv}{good},
  yticklabel style={font=\footnotesize}, enlarge y limits={abs=0.6},
  xlabel={Fall in the quantity bought for a $1\,\%$ rise in the price (percent)},
  nodes near coords, nodes near coords style={font=\scriptsize, anchor=west},
  every node near coord/.append style={/pgf/number format/fixed, /pgf/number format/precision=2,
    /pgf/number format/fixed zerofill}]
\addplot[fill=omDef, draw=none] table[x=elasticity, y expr=\coordindex]
  {figdata/grade-9/02-elasticity-first-look/out/elasticities.csv};
\end{axis}
\end{tikzpicture}
\end{openingchart}

When the \omterm{def:g7:money:price}{price} of beef rises by $1\,\%$, people in the United States buy about $0.75\,\%$ less beef. When
the \omterm{def:g7:money:price}{price} of eggs rises by $1\,\%$, they buy only about $0.27\,\%$ fewer eggs. In the first year after
gasoline becomes $1\,\%$ dearer, drivers buy only about $0.09\,\%$ less; after several years, about
$0.31\,\%$ less. All these buyers respond to the \omterm{def:g7:money:price}{price}, but some much more than others. One number
measures how much.

\begin{recall}
A \omterm{def:g9:demand-and-supply-equations:demand-function}{demand function} gives the \omterm{def:g7:demand-and-supply-first-look:quantity-demanded}{quantity demanded} at each \omterm{def:g7:money:price}{price} (\cref{ch:g9:demand-and-supply-equations}).
A rise of $10\,\%$ multiplies a quantity by $1.10$, a fall of $4\,\%$ by $0.96$. A \omterm{def:g8:taxes-and-public-spending:tax}{tax} on cigarettes lowers
the quantity smoked (\cref{ch:g8:taxes-and-public-spending}).
\end{recall}

\section{How strongly buyers respond}

\begin{definition}[Price elasticity of demand]\label{def:g9:elasticity-first-look:price-elasticity}
The \emph{price elasticity of demand}\index{price elasticity of demand} of a good is the percentage fall
in the \omterm{def:g7:demand-and-supply-first-look:quantity-demanded}{quantity demanded} divided by the percentage rise in its \omterm{def:g7:money:price}{price}, other things being equal:
\[
\text{elasticity} = \frac{\text{percentage fall in the quantity demanded}}{\text{percentage rise in the
price}}.
\]
\end{definition}

If a \omterm{def:g7:money:price}{price} rises by $10\,\%$ and the quantity bought falls by $4\,\%$, the elasticity is $4 \div 10 =
0.4$. Many studies write elasticities with a minus sign ($-0.4$), because \omterm{def:g7:money:price}{price} and quantity move in
opposite directions; this book uses the positive number. In this volume, percentage changes are always
measured from the starting point.

\begin{omfigure}
\includegraphics[width=0.5\linewidth]{images/book1/ai/g9-elasticity-fuel-station.jpg}

{\small A fuel station. When fuel becomes dearer, most drivers still have to reach work the next day:
in the short run, they buy only a little less.}
\end{omfigure}

\section{Elastic and inelastic}

\begin{definition}[Elastic demand]\label{def:g9:elasticity-first-look:elastic}
Demand is \emph{elastic}\index{elastic demand} when its \omterm{def:g7:money:price}{price} elasticity is greater than 1: the quantity
falls by a larger percentage than the \omterm{def:g7:money:price}{price} rises.
\end{definition}

\begin{definition}[Inelastic demand]\label{def:g9:elasticity-first-look:inelastic}
Demand is \emph{inelastic}\index{inelastic demand} when its \omterm{def:g7:money:price}{price} elasticity is smaller than 1: the
quantity falls by a smaller percentage than the \omterm{def:g7:money:price}{price} rises.
\end{definition}

\begin{definition}[Unit-elastic demand]\label{def:g9:elasticity-first-look:unit-elastic}
Demand is \emph{unit-elastic}\index{unit-elastic demand} when its \omterm{def:g7:money:price}{price} elasticity equals 1.
\end{definition}

Every estimate in the opening chart is below 1: the demand for these goods, taken as whole groups, is
\omterm{def:g9:elasticity-first-look:inelastic}{inelastic}. The demand for one brand of a good is usually much more \omterm{def:g9:elasticity-first-look:elastic}{elastic} than the demand for the good
as a whole, because buyers can switch to another brand.

\begin{omfigure}
\begin{tikzpicture}[econ, x=0.5cm, y=0.5cm]
  \begin{scope}
    \draw[econ axes] (0,7) node[above, font=\scriptsize] {price} |- (9,0) node[right, font=\scriptsize] {quantity};
    \draw[demand] (3,7) -- (4.5,0.5);
    \draw[guide] (0,4) -- (3.69,4); \draw[guide] (0,5) -- (3.46,5);
    \draw[guide] (3.69,4) -- (3.69,0); \draw[guide] (3.46,5) -- (3.46,0);
    \node[font=\small\bfseries, anchor=south] at (4.5,7.6) {Inelastic};
    \node[font=\footnotesize, align=center, anchor=north] at (4.5,-0.6) {price up a lot,\\quantity down a little};
  \end{scope}
  \begin{scope}[shift={(12,0)}]
    \draw[econ axes] (0,7) node[above, font=\scriptsize] {price} |- (9,0) node[right, font=\scriptsize] {quantity};
    \draw[demand] (0.5,6) -- (8.5,2.5);
    \draw[guide] (0,4) -- (5.07,4); \draw[guide] (0,5) -- (2.79,5);
    \draw[guide] (5.07,4) -- (5.07,0); \draw[guide] (2.79,5) -- (2.79,0);
    \node[font=\small\bfseries, anchor=south] at (4.5,7.6) {Elastic};
    \node[font=\footnotesize, align=center, anchor=north] at (4.5,-0.6) {the same price rise,\\quantity down a lot};
  \end{scope}
\end{tikzpicture}

{\small The same rise in \omterm{def:g7:money:price}{price} (from the lower to the upper dotted line) along a steep demand line and a
flat one.}
\end{omfigure}

\section{Price and revenue}

A seller's \omterm{def:g7:firms:revenue}{revenue} is the \omterm{def:g7:money:price}{price} times the quantity sold. When the \omterm{def:g7:money:price}{price} rises, each unit brings more
but fewer units are sold. Which effect wins depends on the elasticity. With an elasticity of $0.4$, a
$10\,\%$ \omterm{def:g7:money:price}{price} rise lowers the quantity by $4\,\%$: \omterm{def:g7:firms:revenue}{revenue} is multiplied by $1.10 \times 0.96 = 1.056$,
a rise of $5.6\,\%$. With an elasticity of $1.5$, the quantity falls by $15\,\%$: \omterm{def:g7:firms:revenue}{revenue} is multiplied by
$1.10 \times 0.85 = 0.935$, a fall of $6.5\,\%$.

\begin{model}{Elastic and inelastic demand}
\assumes Demand is stable while the \omterm{def:g7:money:price}{price} changes: other things that affect buyers stay the same.
\predicts A \omterm{def:g7:money:price}{price} rise raises the seller's \omterm{def:g7:firms:revenue}{revenue} when demand is \omterm{def:g9:elasticity-first-look:inelastic}{inelastic} and lowers it when demand
is \omterm{def:g9:elasticity-first-look:elastic}{elastic}.
\tested Cigarettes in South Africa, 1991--2012: as \omterm{def:g8:taxes-and-public-spending:tax}{taxes} and \omterm{def:g7:money:price}{prices} rose, legal sales fell by
$42\,\%$, yet the cigarette makers' \omterm{def:g7:firms:revenue}{revenue} after inflation rose by $95\,\%$ (National Cancer Institute
and World Health Organization, 2016).
\breaks When other things change at the same time as the \omterm{def:g7:money:price}{price} (\omterm{def:g7:households:income}{incomes}, the number of buyers, the
\omterm{def:g7:money:price}{prices} of other goods); when buyers respond more after some years than in the first months.
\end{model}

% figdata: revenue (quantity = 1200 - 100p: p 4 3200, 5 3500, 6 3600, 7 3500, 8 3200)
Along an affine demand line, the elasticity is not the same everywhere: at high \omterm{def:g7:money:price}{prices} demand is \omterm{def:g9:elasticity-first-look:elastic}{elastic},
at low \omterm{def:g7:money:price}{prices} \omterm{def:g9:elasticity-first-look:inelastic}{inelastic}. For the invented demand $Q^d = 1\,200 - 100p$, \omterm{def:g7:firms:revenue}{revenue} rises with the \omterm{def:g7:money:price}{price} up
to 6 dalens, then falls.

\begin{omchart}{Revenue along the invented demand $Q^d = 1\,200 - 100p$: it is largest at a price of 6
dalens, where demand is unit-elastic.}{Invented}
\begin{tikzpicture}
\begin{axis}[omrecord, ybar, bar width=10pt, xmin=-0.6, xmax=12.6, ymin=0, ymax=4000,
  xtick={0,...,12}, xlabel={Price (dalens)}, ylabel={Revenue (dalens)}, height=0.42\linewidth,
  scaled y ticks=false]
\addplot[fill=omDef, draw=none] table[x=price, y=revenue] {figdata/grade-9/02-elasticity-first-look/out/revenue.csv};
\end{axis}
\end{tikzpicture}
\end{omchart}

\begin{measure}[An elasticity from data]
Measuring an elasticity means comparing the quantity bought at two \omterm{def:g7:money:price}{prices}. The difficulty is ``other
things being equal'': between two years, \omterm{def:g7:households:income}{incomes}, tastes and the \omterm{def:g7:money:price}{prices} of other goods change too. Studies
compare places or periods in which the \omterm{def:g7:money:price}{price} changed for a reason unrelated to demand, such as a change
in a \omterm{def:g8:taxes-and-public-spending:tax}{tax}, and check that other things moved little. Different studies of the same good give different
numbers: the foods of the opening chart are averages of 13 to 51 estimates each.
\end{measure}

\section{Necessities and luxuries}

\begin{definition}[Necessity]\label{def:g9:elasticity-first-look:necessity}
A \emph{necessity}\index{necessity} is a good on which \omterm{def:g7:households:household}{households} spend a smaller share of their \omterm{def:g7:households:income}{income}
as their \omterm{def:g7:households:income}{income} rises.
\end{definition}

\begin{definition}[Luxury]\label{def:g9:elasticity-first-look:luxury}
A \emph{luxury}\index{luxury} is a good on which \omterm{def:g7:households:household}{households} spend a larger share of their \omterm{def:g7:households:income}{income} as their
\omterm{def:g7:households:income}{income} rises.
\end{definition}

Food as a whole is a \omterm{def:g9:elasticity-first-look:necessity}{necessity}: richer \omterm{def:g7:households:household}{households} spend a smaller share of their \omterm{def:g7:households:income}{income} on it
(\cref{ch:g7:households}). Goods that \omterm{def:g7:households:household}{households} cannot do
without, and that have no close substitute, tend to have \omterm{def:g9:elasticity-first-look:inelastic}{inelastic demand}; goods that can be postponed or
replaced, like a restaurant meal, more \omterm{def:g9:elasticity-first-look:elastic}{elastic demand}. The opening chart shows the pattern among foods:
meals away from home have the highest elasticity, eggs the lowest.

\section{Short run and long run}

In the first months after a \omterm{def:g7:money:price}{price} rise, buyers can change little: the car, the boiler and the journey to
work are already chosen. Over several years, they can buy a smaller car, insulate the house or move
closer to work. Demand is therefore more \omterm{def:g9:elasticity-first-look:elastic}{elastic} in the long run than in the short run, as the two
gasoline bars of the opening chart show ($0.09$ and $0.31$).

\begin{record}[Fuel prices in 2022]
% ledger: g9:eia-gasoline-2022, g9:eia-gasoline-supplied; figdata: gasoline (2021 3.05, 2022 4.09), facts
% (gas_rise_pct 34.1). EIA's own 2022 average (3.95, in the ledger row) differs from the BLS average printed here
% (two separate price surveys); the BLS series is the one charted.
In the United States, a gallon of regular gasoline cost $3.05$ dollars on average in 2021 and $4.09$ in
2022, a rise of $34.1\,\%$ (Bureau of Labor Statistics average \omterm{def:g7:money:price}{prices}). With a short-run elasticity of
$0.09$, a $34.1\,\%$ rise in the \omterm{def:g7:money:price}{price} predicts a fall in the quantity of only about $3\,\%$: most
drivers keep driving, and pay much more. The fall was smaller still: $8.810$ million barrels of gasoline
a day in 2022 against $8.816$ in 2021, less than $0.1\,\%$ (U.S. Energy Information Administration).
Other things did not stay the same: driving was still coming back after 2020, when consumption had
dropped to $8.049$ million barrels a day, so demand was rising while the \omterm{def:g7:money:price}{price} climbed.
\end{record}

\section{Exercises}

\begin{exercise}[$\star$]\label{exo:g9:elasticity-first-look:1}
Use the opening chart. Which food has the most \omterm{def:g9:elasticity-first-look:elastic}{elastic demand}, and which the least? Which bar of the
whole chart is the lowest?
\end{exercise}

\begin{exercise}[$\star$]\label{exo:g9:elasticity-first-look:2}
The \omterm{def:g7:money:price}{price} of a good rises by $10\,\%$ and the quantity bought falls by $5\,\%$. Compute the \omterm{def:g9:elasticity-first-look:price-elasticity}{price
elasticity of demand}.
\end{exercise}

\begin{exercise}[$\star$]\label{exo:g9:elasticity-first-look:3}
Is demand \omterm{def:g9:elasticity-first-look:elastic}{elastic} or \omterm{def:g9:elasticity-first-look:inelastic}{inelastic} when the elasticity is $0.5$? When it is $1.8$?
\end{exercise}

\begin{exercise}[$\star$]\label{exo:g9:elasticity-first-look:4}
A shop sells 500 units at 3 dalens, or 400 units at 4 dalens. Compute its \omterm{def:g7:firms:revenue}{revenue} at each \omterm{def:g7:money:price}{price}.
\end{exercise}

\begin{exercise}[$\star\star$]\label{exo:g9:elasticity-first-look:5}
A \omterm{def:g7:money:price}{price} rises by $10\,\%$; the elasticity is $0.4$. By what percentage do the quantity and the \omterm{def:g7:firms:revenue}{revenue}
change?
\end{exercise}

\begin{exercise}[$\star\star$]\label{exo:g9:elasticity-first-look:6}
The same, with an elasticity of $1.5$.
\end{exercise}

\begin{exercise}[$\star\star$]\label{exo:g9:elasticity-first-look:7}
When the \omterm{def:g7:money:price}{price} goes from 4 to 5 dalens, the quantity bought goes from 800 to 700 units. Compute the
elasticity, measuring both changes from the starting point.
\end{exercise}

\begin{exercise}[$\star\star$]\label{exo:g9:elasticity-first-look:8}
Use the gasoline elasticities of the opening chart. By about what percentage does a $10\,\%$ rise in the
\omterm{def:g7:money:price}{price} of gasoline lower the quantity bought in the first year? After several years?
\end{exercise}

\begin{exercise}[$\star\star$]\label{exo:g9:elasticity-first-look:9}
A city raises its bus fares by $10\,\%$. The same year, a large road bridge closes for repairs, and the
number of bus rides rises by $5\,\%$. Does this show that the demand for bus rides does not fall when the
fare rises?
\end{exercise}

\begin{exercise}[$\star\star\star$]\label{exo:g9:elasticity-first-look:10}
For the demand $Q^d = 1\,200 - 100p$, compute the \omterm{def:g7:firms:revenue}{revenue} at \omterm{def:g7:money:price}{prices} of 4, 5, 6, 7 and 8 dalens. Which
\omterm{def:g7:money:price}{price} gives the largest \omterm{def:g7:firms:revenue}{revenue}?
\end{exercise}

\begin{exercise}[$\star\star\star$]\label{exo:g9:elasticity-first-look:11}
Two goods each sell $1\,000$ units a year at 10 dalens. A \omterm{def:g8:taxes-and-public-spending:tax}{tax} of 1 dalen a unit raises each \omterm{def:g7:money:price}{price} by
$10\,\%$. The elasticity is $0.3$ for the first good and $2.0$ for the second. How much \omterm{def:g8:taxes-and-public-spending:tax}{tax} does each raise
in a year?
\end{exercise}

\begin{exercise}[$\star\star\star$]\label{exo:g9:elasticity-first-look:12}
Give two reasons why an elasticity measured by comparing two years may mislead.
\end{exercise}

\section{Problem: The Bus Fare of Brenmoor}

\begin{problem}[{Weekend problem --- riders at two fares, the elasticity and revenue, and a fare cut in
the short and the long run}]
\label{pb:g9:elasticity-first-look:1}
Everything in this problem is invented. Brenmoor's buses carry $10\,000$ riders a day at a fare of
$2.00$ dalens. The town raises the fare to $2.50$ dalens; the next month, $9\,000$ riders a day take the
bus.

\medskip
\noindent\textbf{Part I --- Two fares.}
\begin{enumerate}
  \item Compute the daily \omterm{def:g7:firms:revenue}{revenue} at each fare.
  \item By what percentage did the fare rise?
  \item By what percentage did the number of riders fall?
  \item By how much did the daily \omterm{def:g7:firms:revenue}{revenue} change?
\end{enumerate}

\medskip
\noindent\textbf{Part II --- Elasticity and revenue.}
\begin{enumerate}[resume]
  \item Compute the elasticity of the demand for bus rides.
  \item Is this demand \omterm{def:g9:elasticity-first-look:elastic}{elastic} or \omterm{def:g9:elasticity-first-look:inelastic}{inelastic}?
  \item If the elasticity had been $1.5$, how many riders would there have been at $2.50$ dalens, and what
        \omterm{def:g7:firms:revenue}{revenue}?
  \item Explain why the \omterm{def:g7:firms:revenue}{revenue} rose in question 4 but falls in question 7.
\end{enumerate}

\medskip
\noindent\textbf{Part III --- A fare cut.} A year later, the town cuts the fare from $2.50$ back to $2.00$
dalens.
\begin{enumerate}[resume]
  \item With the short-run elasticity of question 5, measured from $2.50$ dalens, how many riders are there
        after the cut? What is the daily \omterm{def:g7:firms:revenue}{revenue}?
  \item By how much does the daily \omterm{def:g7:firms:revenue}{revenue} change in the short run?
  \item Over several years, people sell cars and move along the bus lines: the elasticity becomes $1.2$.
        How many riders are there then, and what is the daily \omterm{def:g7:firms:revenue}{revenue}?
  \item By how much has the daily \omterm{def:g7:firms:revenue}{revenue} changed in the long run, compared with the \omterm{def:g7:firms:revenue}{revenue} at $2.50$
        dalens?
\end{enumerate}
\end{problem}
